When a language model answers a factual question, where does its internal state go? We asked four open models (Llama-3.1-8B, Qwen2.5-7B, Gemma-2-9B and Llama-3.2-1B) \NFacts{} factual questions whose answers recur, such as the country a person is a citizen of. At every layer, we recorded the state from which the answer is produced. Questions that receive the same answer converge on a shared region of state space, which we summarise by its centre, the \emph{answer centroid}. We report four findings. First, the answer becomes readable in a single layer. In every model, the share of questions whose state lies nearest its own answer's centroid jumps within one layer, at a depth between \OnsetMin{} and \OnsetMax{} of the way through the network. Second, centroids follow the model's answer, not the question: the centroid of a country learned from questions about people's citizenship also identifies that country in questions about towns and films. Third, the model uses the centroids. After the jump, moving the state of a wrongly answered question to the right answer's centroid makes the model give the right answer in \PushGoldMin{}--\PushGoldMax{} of cases late in the network. A random move of the same size does so in \PushRandomMin{}--\PushRandomMax{}. Before the jump, moves of the same size almost never give the right answer. Fourth, near misses are visible from the onset. When a model answers wrongly but the right answer's centroid is among the nearest to its state, the model usually picks the right answer in multiple choice. This signal, read partway through the network, is about as informative as the right answer's rank in the model's final output, and more informative than the model's confidence in its reply.
